\section{Introduction}
In recent years, the frontier of Large Language Models (LLMs) has shifted decisively toward enhancing their reasoning capabilities, pushing the boundaries of Artificial General Intelligence (AGI). 
State-of-the-art models such as OpenAI's o1~\citep{openai_o1}, OpenAI's o3 \citep{o3}, Gemini 2.5 \citep{comanici2025gemini}, DeepSeek-R1 \citep{guo2025deepseek}, Qwen3 \citep{yang2025qwen3}, and GLM-4.5 \citep{5team2025glm45agenticreasoningcoding} showcase this trend, demonstrating profound abilities in complex logic, mathematics, code generation, and agentic tasks. This advancement is largely driven by a new paradigm: leveraging large-scale Reinforcement Learning (RL)~\citep{sutton1998reinforcement} not only to refine models but also to enable deeper, more extensive reasoning at inference time. By dynamically allocating more computational FLOPs to extend Chain-of-Thought (CoT)~\citep{wei2022chain}, specialized models like OpenAI's o1 and Gemini 2.5 Pro have achieved breakthrough performance on formidable challenges, including olympiad-level mathematics, competitive programming, and complex agentic scenarios.

In this work, we introduce {\longcatthink}, an efficient open-source Mixture-of-Experts (MoE) reasoning model that establishes a new state-of-the-art for open-source reasoning models. Built upon our foundational {\longcatbase} model~\citep{meituan2025longcat_flash_chat}, {\longcatthink} ($560$B total parameters: 27B active on average) excels on complex logic, STEM, coding, and agent tasks. The development of {\longcatthink} follows a meticulously designed two-phase pipeline: Long CoT Cold-Start Training and Large-Scale RL. In the first phase, we aim to build the model's foundational reasoning abilities. This begins with a curriculum learning strategy during mid-training to strengthen intrinsic capabilities, followed by a targeted Supervised Fine-Tuning (SFT) stage on reasoning-intensive and agentic data to prepare the model for advanced learning. The second phase scales up this potential through an efficient RL framework, built upon our Dynamic ORchestration for Asynchronous rollout (DORA) system for industrial-scale asynchronous training. To address the stability challenges in asynchronous RL training, we adapt and extend the Group Relative Policy Optimization (GRPO)~\citep{grpo} algorithm for a robust exploration-exploitation balance. A key innovation in this phase is our domain-parallel training scheme, which simultaneously optimizes the model across distinct domains and subsequently merges the resulting domain-expert models into a fused model~\citep{yu2024language}. A final general RL stage further refines the fused model, improving its robustness, safety, and human alignment for real-world applications. The overview of the training pipeline is illustrated in Figure~\ref{fig:parallel_training_pipeline}.

\begin{figure}[h]
\centerline{\includegraphics[width=\linewidth]{figs/parallel_training_pipeline}}
% \vspace{-.25cm}
\caption{\small
The training pipeline of~{\longcatthink}. We first perform mid-training and SFT to cultivate the base model’s foundational reasoning abilities. Then, we introduce a domain-parallel training scheme during the large-scale RL phase to obtain multiple domain-specific experts.
Finally, a general RL stage followed by expert models fusion to improve the general ability.
}
\label{fig:parallel_training_pipeline}
\end{figure}

Our work presents three core contributions:
\begin{itemize}
    \item \textbf{Domain-Parallel RL Training and Fusion Methodology}: To overcome the instability of traditional mixed-domain RL training, we design a domain-parallel scheme that decouples optimization across STEM, coding, and agentic tasks. This approach not only stabilizes training but also allows us to fuse the resulting domain-expert models into a nearly Pareto-optimal final model that excels across all specialties.
    \item \textbf{Pioneering Industrial-Scale RL Infrastructure}:  Our DORA system provides the robust backbone for RL training. Its asynchronous architecture delivers a greater than threefold speedup over synchronous frameworks, enabling stable training on tens of thousands of accelerators. This industrial-grade system supported a massive RL investment of nearly $20\%$ of pre-training compute, making our advanced methodology feasible at scale. We also highlight its novel features, including elastic colocation and efficient KV-cache reuse, which are detailed in Section~\ref{subsec:rl_infra}.
    \item \textbf{Broad and Efficient Advanced Reasoning}: We significantly extend the model's capabilities into challenging domains, achieving exceptional performance and efficiency. To improve agentic capabilities, we propose a dual-path reasoning approach to select high-value training queries that benefit most from tool integration. We complement this with an automated pipeline to construct high-quality, tool-augmented reasoning trajectories for training. For agentic tool use, {\longcatthink} leads to a $64.5\%$ reduction in token consumption on the AIME-25 benchmark while preserving accuracy. 
    For formal reasoning, we developed an expert-iteration pipeline integrated with a Lean4 server to synthetically generate verified proofs, systematically instilling a capability absent in most LLMs.
\end{itemize}
These contributions culminate in a model that not only achieves state-of-the-art performance on a diverse set of benchmarks (Figure~\ref{fig:qualitative_example}), but also establishes a clear advantage in the crucial domains of formal proving and agentic reasoning over its open-source peers.
